% GENERATED FILE. Edit canonical lessons, then run npm run generate:books.
% canonical-source-hash: fnv1a64:ebde5244f0f6923b
% canonical-lessons: TA-C164-vivari, TA-C164-oppatai, TA-C164-payirci-cey, TA-C164-tercci-peru, TA-C164-kavani

\chapter{To describe, To hand over, To practise, To pass, To pay attention}
\label{ch:ta-to-describe-to-hand-over-to-practise-to-pass-to-pay-attention}
\hlchaptermodality{164}

\begin{chapteropening}
\textbf{By the end of this chapter:} \emph{I can say to describe, to hand over, to practise, to pass and to pay attention.}

Complete the last lesson of chapter 164: I can say to describe, to hand over, to practise, to pass and to pay attention.
\end{chapteropening}

\section[\texorpdfstring{\ta{விவரி} (vivari)}{விவரி (vivari)}]{\ta{விவரி} (vivari) — to describe, to explain}

\label{lesson:TA-C164-vivari}



\begin{quote}
Before the new one: say the Tamil for to bargain, then the Tamil for to check.
\end{quote}

\subsection*{You'll want to know: \ta{விவரி}}
\textbf{\ta{விவரி}} — \emph{vivari} — \textquotedblleft{}to describe, to explain\textquotedblright{}.

\begin{grammarlens}[title={What you've built}]
One more word for this chapter.
\end{grammarlens}

\subsection*{Your turn}
\begin{itemize}
\raggedright
  \item \emph{Say it:} \emph{vivari}
  \item \emph{Say it:} \emph{vivari}, once more
  \item \emph{Say it:} say \emph{vivari}
  \item \emph{Recall:} say the Tamil for to bargain, then the Tamil for to check, then say \emph{vivari} again
\end{itemize}

\begin{tcolorbox}[breakable,colback=teal!4,colframe=teal!35!black,title={Before you move on}]
What does \ta{விவரி} mean? (\textquotedblleft{}To describe, to explain\textquotedblright{}.) Say it once more.
\end{tcolorbox}
\section[\texorpdfstring{\ta{ஒப்படை} (oppaṭai)}{ஒப்படை (oppaṭai)}]{\ta{ஒப்படை} (oppaṭai) — to hand over, to hand in}

\label{lesson:TA-C164-oppatai}



\begin{quote}
Before the new one: say the Tamil for to check, then the Tamil for to describe.
\end{quote}

\subsection*{You'll want to know: \ta{ஒப்படை}}
\textbf{\ta{ஒப்படை}} — \emph{oppaṭai} — \textquotedblleft{}to hand over, to hand in\textquotedblright{}.

\begin{grammarlens}[title={What you've built}]
One more word for this chapter.
\end{grammarlens}

\subsection*{Your turn}
\begin{itemize}
\raggedright
  \item \emph{Say it:} \emph{oppaṭai}
  \item \emph{Say it:} \emph{oppaṭai}, once more
  \item \emph{Say it:} say \emph{oppaṭai}
  \item \emph{Recall:} say the Tamil for to check, then the Tamil for to describe, then say \emph{oppaṭai} again
\end{itemize}

\begin{tcolorbox}[breakable,colback=teal!4,colframe=teal!35!black,title={Before you move on}]
What does \ta{ஒப்படை} mean? (\textquotedblleft{}To hand over, to hand in\textquotedblright{}.) Say it once more.
\end{tcolorbox}
\section[\texorpdfstring{\ta{பயிற்சி} \ta{செய்} (payiṟci cey)}{பயிற்சி செய் (payiṟci cey)}]{\ta{பயிற்சி} \ta{செய்} (payiṟci cey) — to practise}

\label{lesson:TA-C164-payirci-cey}



\begin{quote}
Before the new one: say the Tamil for to describe, then the Tamil for to hand over.
\end{quote}

\subsection*{You'll want to know: \ta{பயிற்சி} \ta{செய்}}
\textbf{\ta{பயிற்சி} \ta{செய்}} — \emph{payiṟci cey} — \textquotedblleft{}to practise\textquotedblright{}.

\begin{grammarlens}[title={What you've built}]
One more word for this chapter.
\end{grammarlens}

\subsection*{Your turn}
\begin{itemize}
\raggedright
  \item \emph{Say it:} \emph{payiṟci cey}
  \item \emph{Say it:} \emph{payiṟci cey}, once more
  \item \emph{Say it:} say \emph{payiṟci cey}
  \item \emph{Recall:} say the Tamil for to describe, then the Tamil for to hand over, then say \emph{payiṟci cey} again
\end{itemize}

\begin{tcolorbox}[breakable,colback=teal!4,colframe=teal!35!black,title={Before you move on}]
What does \ta{பயிற்சி} \ta{செய்} mean? (\textquotedblleft{}To practise\textquotedblright{}.) Say it once more.
\end{tcolorbox}
\section[\texorpdfstring{\ta{தேர்ச்சி} \ta{பெறு} (tērcci peṟu)}{தேர்ச்சி பெறு (tērcci peṟu)}]{\ta{தேர்ச்சி} \ta{பெறு} (tērcci peṟu) — to pass (an exam)}

\label{lesson:TA-C164-tercci-peru}



\begin{quote}
Before the new one: say the Tamil for to hand over, then the Tamil for to practise.
\end{quote}

\subsection*{You'll want to know: \ta{தேர்ச்சி} \ta{பெறு}}
\textbf{\ta{தேர்ச்சி} \ta{பெறு}} — \emph{tērcci peṟu} — \textquotedblleft{}to pass (an exam)\textquotedblright{}.

\begin{grammarlens}[title={What you've built}]
One more word for this chapter.
\end{grammarlens}

\subsection*{Your turn}
\begin{itemize}
\raggedright
  \item \emph{Say it:} \emph{tērcci peṟu}
  \item \emph{Say it:} \emph{tērcci peṟu}, once more
  \item \emph{Say it:} say \emph{tērcci peṟu}
  \item \emph{Recall:} say the Tamil for to hand over, then the Tamil for to practise, then say \emph{tērcci peṟu} again
\end{itemize}

\begin{tcolorbox}[breakable,colback=teal!4,colframe=teal!35!black,title={Before you move on}]
What does \ta{தேர்ச்சி} \ta{பெறு} mean? (\textquotedblleft{}To pass (an exam)\textquotedblright{}.) Say it once more.
\end{tcolorbox}
\section[\texorpdfstring{\ta{கவனி} (kavaṉi)}{கவனி (kavaṉi)}]{\ta{கவனி} (kavaṉi) — to pay attention, to notice}

\label{lesson:TA-C164-kavani}



\begin{quote}
Before the new one: say the Tamil for to practise, then the Tamil for to pass.
\end{quote}

\subsection*{You'll want to know: \ta{கவனி}}
\textbf{\ta{கவனி}} — \emph{kavaṉi} — \textquotedblleft{}to pay attention, to notice\textquotedblright{}.

\begin{grammarlens}[title={What you've built}]
One more word for this chapter.
\end{grammarlens}

\subsection*{Your turn}
\begin{itemize}
\raggedright
  \item \emph{Say it:} \emph{kavaṉi}
  \item \emph{Say it:} \emph{kavaṉi}, once more
  \item \emph{Say it:} say \emph{kavaṉi}
  \item \emph{Recall:} say the Tamil for to practise, then the Tamil for to pass, then say \emph{kavaṉi} again
\end{itemize}

\begin{tcolorbox}[breakable,colback=teal!4,colframe=teal!35!black,title={Before you move on}]
What does \ta{கவனி} mean? (\textquotedblleft{}To pay attention, to notice\textquotedblright{}.) Say it once more.
\end{tcolorbox}
